% Generated draft: requires mathematical and source review.
\documentclass[11pt]{article}
\usepackage[margin=25mm]{geometry}
\usepackage{fontspec}
\setmainfont{Latin Modern Roman}
\usepackage{amsmath,amssymb,mathtools,unicode-math}
\setmathfont{Latin Modern Math}
\usepackage{enumitem,needspace,xurl,hyperref}
\setlength{\parindent}{0pt}
\setlength{\parskip}{5pt}
\setlength{\emergencystretch}{3em}
\newcommand{\N}{\mathbb{N}}
\newcommand{\Z}{\mathbb{Z}}
\newcommand{\Q}{\mathbb{Q}}
\newcommand{\R}{\mathbb{R}}
\newcommand{\C}{\mathbb{C}}
\newcommand{\F}{\mathbb{F}}
\newcommand{\A}{\mathbb{A}}
\newcommand{\PP}{\mathbb{P}}
\newcommand{\E}{\mathbb{E}}
\newcommand{\eps}{\varepsilon}
\newcommand{\dd}{\,\mathrm{d}}
\begin{document}
\section*{UnsolvedMath ID: 30000573}
\textbf{Source problem code:} catalogue entry
\subsection{Definitions and mathematical statement}
\textbf{Chaotic-operator question.} Does every infinite-dimensional nuclear Fr\'echet space admit a chaotic operator? A Fr\'echet space \(E\) is a complete metrizable Hausdorff locally convex vector space over the intended real or complex scalar field. In the usual seminorm criterion, nuclearity means that for every continuous seminorm \(p\) there is a continuous seminorm \(q\ge p\) such that the canonical map \(E_q\to E_p\) between the corresponding Banach completions of seminorm quotients is nuclear. A nuclear Banach-space map admits a series \(\sum_j\lambda_j\ell_j(\cdot)y_j\) with \(\sum_j|\lambda_j|<\infty\) and uniformly bounded functionals and vectors.\par Under the standard linear-dynamics meaning of chaotic, the requested operator is a continuous linear map \(T:E\to E\) satisfying both\[\exists x\in E:\quad\overline{\{T^nx:n\in\Z_{\ge0}\}}=E,\qquad \overline{\bigcup_{m\ge1}\ker(T^m-I)}=E.\] The first property is hypercyclicity and the second is density of periodic points. Closures use the original Fr\'echet topology. Nuclearity is a property of \(E\), not a requirement that \(T\) itself be a nuclear operator. No Schauder basis, continuous norm, or special sequence-space representation is supplied or imposed. A hypercyclic operator without dense periodic points would not answer the stated chaotic-operator question. The source does not fix the scalar field or explicitly define chaos, so those standard conventions need source confirmation.\par The quantifiers are \(\forall E\) in the stated class, \(\exists T\) with both properties; the operator may depend on the space. The supplied triage checked 2026-08-21 labels the problem open and reports no located authoritative resolution, using the report alone as its reference. No precise later construction or obstruction is supplied. The dated status claim is preserved without independently asserting that the question remains unresolved today.
\subsection{Short English statement}
On each infinite-dimensional nuclear Fr\'echet space, seek a continuous linear operator with both a dense orbit and dense periodic points. Hypercyclicity alone is insufficient, and no extra basis or norm hypothesis is inserted. The supplied open-status assessment is dated and unverified here.
\subsection{Sources}
\begin{itemize}\item Local UnsolvedMath registry, permanent ID 30000573, code catalogue entry; \emph{Chaotic Operators on Nuclear Fr\'echet Spaces}, Oberwolfach Reports Open Problems.\item \emph{Mini-Workshop: Hypercyclicity and Linear Chaos (2007)}, Oberwolfach Reports, DOI 10.4171/owr/2006/37, pp. 2227--2276. \url{https://doi.org/10.4171/owr/2006/37}. Proposed year and year tag: 2006; no proposer named.\item The same full report is the sole supplied literature source; no separate authoritative chaotic-operator theorem is supplied.\item Literature triage checked 2026-08-21, status \texttt{open}; record updated 2026-09-16. Statement status \texttt{exact}, reported retained byte-for-byte. Verify dated claims independently.\end{itemize}
\end{document}
